% Ejercicios nuevos para el 2do Parcial — Química General A1095 (UTN FRBA, Ing. Mecánica).
% Mismos temas que el repaso de cátedra, con otros datos; respuestas en la última página.
% Compilar con: make  ->  build/ejercicios_nuevos.pdf

\documentclass[11pt]{article}
\usepackage[a4paper, margin=20mm]{geometry}
\usepackage{amsmath}
\usepackage{fontspec}
\usepackage{unicode-math}
\setmainfont{Latin Modern Roman}
\setmathfont{Latin Modern Math}
\usepackage[dvipsnames]{xcolor}
\usepackage{enumitem}
\usepackage[most]{tcolorbox}
\usepackage{array}

\pagestyle{empty}
\frenchspacing
\setlength{\parindent}{0pt}
\setlength{\parskip}{3pt}
\setlist[enumerate,1]{label=\textbf{\alph*)}, leftmargin=2.2em, itemsep=2pt, topsep=3pt}
\setlist[enumerate,2]{label=\textbf{\roman*)}, leftmargin=2em, itemsep=0pt, topsep=1pt}

% Ejercicio: título con filete.
\newcommand{\ejercicio}[1]{\par\vspace{10pt}{\large\bfseries\color{NavyBlue}Ejercicio #1}\par
  \nointerlineskip\vspace{2pt}{\color{NavyBlue!40}\rule{\linewidth}{0.4pt}}\par\vspace{2pt}}
% Recuadro de datos: fondo gris con barra a la izquierda.
\newtcolorbox{datos}{enhanced, colback=gray!12, colframe=gray!12, boxrule=0pt,
  sharp corners, left=6pt, right=4pt, top=2pt, bottom=2pt, before skip=6pt, after skip=2pt,
  borderline west={2.5pt}{0pt}{black!80}}

\begin{document}

{\LARGE\bfseries Química General — Ejercicios nuevos}\hfill{\small UTN FRBA · Ing. Mecánica}\par
\vspace{2pt}{\color{NavyBlue}\rule{\linewidth}{0.8pt}}\par

\ejercicio{1}
Se tienen 250 g de vapor de agua a \(100{,}0\) °C y 1 atm, que se condensan y luego se enfrían hasta agua líquida a \(40{,}0\) °C.
\begin{enumerate}
\item Calcule el calor total (en kJ) que libera el sistema.
\item ¿Qué masa de hielo a 0 °C se podría fundir con ese calor?
\item ¿Qué masa de butano (C\(_4\)H\(_{10}\)) habría que quemar para obtener la misma cantidad de calor, si \(\Delta H^\circ_c = -2877\) kJ/mol?
\end{enumerate}
\begin{datos}
\textbf{Datos:} \(s_{\text{agua (l)}} = 4{,}18\) J/(g·°C)\,; \(\Delta H^\circ_{\text{fus}} = 334\) J/g\,; \(\Delta H^\circ_{\text{vap}} = 2260\) J/g\,; \(M_r(\text{C}_4\text{H}_{10}) = 58{,}1\) g/mol
\end{datos}

\ejercicio{2}
En el alto horno, el óxido férrico se reduce con monóxido de carbono:
\[\text{Fe}_2\text{O}_3(s) + 3\,\text{CO}(g) \rightarrow 2\,\text{Fe}(s) + 3\,\text{CO}_2(g)\]
\begin{enumerate}
\item Calcule \(\Delta H^\circ_r\) a partir de las entalpías estándar de formación:\par
{\raggedright\(\Delta H^\circ_f[\text{Fe}_2\text{O}_3(s)] = -822{,}2\) kJ/mol\,; \(\Delta H^\circ_f[\text{CO}(g)] = -110{,}5\) kJ/mol\,;\\ \(\Delta H^\circ_f[\text{CO}_2(g)] = -393{,}5\) kJ/mol\par}
\item ¿Por qué no hace falta el dato de \(\Delta H^\circ_f\) del Fe(s)? Indique si el proceso es endotérmico o exotérmico.
\item Calcule la energía (en kJ) liberada o absorbida al producir 1,00 kg de Fe.
\end{enumerate}
\begin{datos}
\textbf{Datos:} \(A_r(\text{Fe}) = 55{,}85\) g/mol
\end{datos}

\ejercicio{3}
En un recipiente rígido de 2,00 L a 350 K se introducen 0,500 moles de N\(_2\)O\(_4\)(g). Al alcanzarse el equilibrio:
\[\text{N}_2\text{O}_4(g) \rightleftharpoons 2\,\text{NO}_2(g) \qquad (\Delta H^\circ > 0)\]
hay 0,300 moles de NO\(_2\)(g) en el recipiente.
\begin{enumerate}
\item Calcule las concentraciones molares de las dos especies en el equilibrio.
\item Calcule \(K_c\) y \(K_p\) a 350 K.
\item En otro recipiente a 350 K hay \([\text{N}_2\text{O}_4] = 0{,}10\) M y \([\text{NO}_2] = 0{,}20\) M. Calcule \(Q_c\) e indique hacia dónde evoluciona el sistema.
\item Prediga el desplazamiento del equilibrio y si cambia \(K\) cuando:
  \begin{enumerate}
  \item se disminuye la temperatura a volumen constante;
  \item se aumenta el volumen del recipiente a temperatura constante;
  \item se agrega helio a volumen constante.
  \end{enumerate}
\end{enumerate}

\newpage
\ejercicio{4}
Se dispone de una solución acuosa 0,15 M de amoníaco (NH\(_3\)) a 25 °C.
\begin{enumerate}
\item Escriba la ecuación de ionización del amoníaco en agua.
\item Calcule \([\text{OH}^-]\), el pOH, el pH y el \% ionizado.
\item Calcule la \(K_a\) del ion amonio (NH\(_4^+\)), su ácido conjugado.
\end{enumerate}
\begin{datos}
\textbf{Datos:} \(K_b(\text{NH}_3) = 1{,}8\times10^{-5}\)\,; \(K_w = 1{,}0\times10^{-14}\)
\end{datos}

\ejercicio{5}
Se mezclan 50,0 mL de H\(_2\)SO\(_4\) 0,100 M con 150 mL de NaOH 0,100 M, y la mezcla se diluye con agua destilada hasta 250 mL.
\begin{enumerate}
\item Escriba la ecuación balanceada de neutralización.
\item Determine qué reactivo está en exceso y cuántos moles sobran.
\item Calcule la concentración del reactivo sobrante en la solución final y el pH (asuma disociación total del H\(_2\)SO\(_4\), con 2 H\(^+\) por fórmula, y volúmenes aditivos).
\item ¿Qué volumen de NaOH 0,100 M habría hecho falta para neutralizar exactamente los 50,0 mL del ácido?
\end{enumerate}

\ejercicio{6}
A partir de los siguientes potenciales estándar de reducción a 25 °C:
\[E^\circ_{\text{red}}(\text{Al}^{3+}/\text{Al}) = -1{,}66\ \text{V} \qquad ; \qquad E^\circ_{\text{red}}(\text{Cu}^{2+}/\text{Cu}) = +0{,}34\ \text{V}\]
\begin{enumerate}
\item Indique cuál es el ánodo y cuál el cátodo, y cuál especie es el mejor agente oxidante.
\item Escriba las hemirreacciones y la reacción global balanceada (ojo: los electrones intercambiados no son iguales en las dos hemirreacciones).
\item Calcule la FEM estándar e indique si la reacción es espontánea. ¿Cambia \(E^\circ\) al multiplicar las hemirreacciones para balancear?
\item Escriba la notación abreviada de la pila.
\end{enumerate}

\ejercicio{7}
Se quiere depositar 5,00 g de níquel en el cátodo por electrólisis de una solución acuosa de NiCl\(_2\), con una corriente continua de 3,00 A.
\begin{enumerate}
\item Escriba las hemirreacciones en el cátodo (Ni) y en el ánodo (Cl\(_2\)).
\item Calcule la carga \(Q\) necesaria y el tiempo (en minutos) que debe circular la corriente.
\item Calcule el volumen de Cl\(_2\) liberado en el ánodo, medido en CNPT (\(V_m = 22{,}4\) L/mol).
\end{enumerate}
\begin{datos}
\textbf{Datos:} \(F = 96\,500\) C/mol e\(^-\)\,; \(A_r(\text{Ni}) = 58{,}7\) g/mol
\end{datos}

\newpage
{\Large\bfseries Respuestas}\hfill{\small\color{gray}solo resultados finales, para verificar}\par
\vspace{2pt}{\color{NavyBlue}\rule{\linewidth}{0.8pt}}\par\vspace{6pt}
\renewcommand{\arraystretch}{1.5}
\begin{tabular}{@{}>{\bfseries\color{NavyBlue}}l >{\raggedright\arraybackslash}p{0.92\linewidth}@{}}
1 & a) \(q = 565{,}0 + 62{,}7 = 627{,}7\) kJ liberados \(\approx 628\) kJ \quad b) 1,88 kg de hielo \quad c) 0,218 mol \(\rightarrow\) 12,7 g de C\(_4\)H\(_{10}\)\\
2 & a) \(\Delta H^\circ_r = -26{,}8\) kJ \quad b) Fe(s) es un elemento en su forma estable: \(\Delta H^\circ_f = 0\); exotérmico \quad c) 17,9 mol de Fe \(\rightarrow\) 240 kJ liberados\\
3 & a) \([\text{N}_2\text{O}_4] = 0{,}175\) M; \([\text{NO}_2] = 0{,}150\) M \quad b) \(K_c = 0{,}129\); \(K_p = 3{,}69\) (\(\Delta n = 1\)) \quad c) \(Q_c = 0{,}40 > K_c\) \(\rightarrow\) hacia reactivos \quad d) i) hacia reactivos, \(K\) baja; ii) hacia productos (más moles de gas), \(K\) no cambia; iii) no se desplaza\\
4 & a) \(\text{NH}_3 + \text{H}_2\text{O} \rightleftharpoons \text{NH}_4^+ + \text{OH}^-\) \quad b) \([\text{OH}^-] = 1{,}6\times10^{-3}\) M; pOH = 2,78; pH = 11,22; 1,1~\% \quad c) \(K_a = 5{,}6\times10^{-10}\)\\
5 & a) \(\text{H}_2\text{SO}_4 + 2\,\text{NaOH} \rightarrow \text{Na}_2\text{SO}_4 + 2\,\text{H}_2\text{O}\) \quad b) NaOH en exceso: 0,0100 mol H\(^+\) vs.\ 0,0150 mol OH\(^-\) \(\rightarrow\) sobran 0,0050 mol OH\(^-\) \quad c) \([\text{OH}^-] = 0{,}020\) M; pH = 12,30 \quad d) 100 mL\\
6 & a) Ánodo Al, cátodo Cu; mejor oxidante Cu\(^{2+}\) \quad b) \(2\,\text{Al} + 3\,\text{Cu}^{2+} \rightarrow 2\,\text{Al}^{3+} + 3\,\text{Cu}\) (6 e\(^-\)) \quad c) FEM = 2,00 V, espontánea; \(E^\circ\) no cambia \quad d) Al(s)\,|\,Al\(^{3+}\)(ac)\,||\,Cu\(^{2+}\)(ac)\,|\,Cu(s)\\
7 & a) \(\text{Ni}^{2+} + 2\,\text{e}^- \rightarrow \text{Ni}\); \(2\,\text{Cl}^- \rightarrow \text{Cl}_2 + 2\,\text{e}^-\) \quad b) \(Q = 1{,}64\times10^4\) C; \(t = 5480\) s \(\approx 91{,}3\) min \quad c) 1,91 L de Cl\(_2\)\\
\end{tabular}

\end{document}
